\section{Riffaut's conjecture and sparse equations for singular moduli}

A singular modulus is a value $\alpha=j(\tau)$ of the modular invariant at
an imaginary quadratic argument in the upper half-plane. The corresponding
elliptic curve has complex multiplication by an order of discriminant
$\Delta<0$. The conjugates of $\alpha$ over $\mathbb Q$ are exactly the
singular moduli of discriminant $\Delta$, and their number is the order's
class number $h(\Delta)$. A polynomial's support size is its number of
nonzero coefficients. Write $s(\Delta)$ for the minimum support size of a
nonzero rational polynomial vanishing at one, and hence every, such modulus.

\begin{theorem}[Transfer to Riffaut's conjecture]
Accepting the Dirichlet GRH theorem exported by the repository,
$s(\Delta)\ge4$ whenever $h(\Delta)\ge3$.
\end{theorem}
\begin{proof}
Suppose $\alpha$ satisfies a relation of support at most three.
It is nonzero since it has degree at least three. Dividing by its smallest
power of $\alpha$ and its leading coefficient gives
$\alpha^u+A\alpha^v+B=0$, with $u>v>0$, $A,B\in\mathbb Q$ and
$B\ne0$. A binomial relation of degree at least three permits $A=0$.

Bilu, Luca and Pizarro-Madariaga call such a $\Delta$ a trinomial
discriminant. Their Theorems 1.3 and 1.4 imply
$m=|\Delta|>10^{11}$ and $\Delta=-p$ or $-pq$ for distinct odd primes
in the second case, with $\Delta$ fundamental. Hence the Kronecker character
$\chi=(\Delta/\cdot)$ is primitive, real, and odd, of conductor $m$ with at
most two distinct prime factors. Their Corollary 5.2 forces every prime $p$
with $\chi(p)=1$ to satisfy $p>4\sqrt m/\log m$. Their Proposition 9.2,
under Dirichlet GRH, produces a prime with $\chi(p)=1$ and
$p\le4\sqrt m/\log m$ for this very character, since $m\ge10^{10}$.
This contradiction is their Theorem 9.1. Its zero-location input is supplied
by the upstream theorems for the Riemann zeta-function and primitive Dirichlet
characters. The results about singular moduli and the explicit constants
are the published authors' work.
\end{proof}

The two arguments meet at a specific arithmetic object. Section 9.1 of
the published proof uses the smoothed sum
$T(x,\chi)=\sum_{n\le x}\chi(n)\Lambda(n)(1/n-1/x)$.
The repository's prime remainder uses the same kernel with weight
$1-\chi_K(n)$ in the progression $1\bmod q$. Its eventual estimate has
parameters depending on the modulus. The explicit conductor threshold in
Proposition 9.2 enters through the published argument after applying GRH.
Full details appear in Y.~Bilu, F.~Luca and A.~Pizarro-Madariaga,
\emph{Trinomials, singular moduli and Riffaut's conjecture},
Ann. Sc. Norm. Super. Pisa Cl. Sci. (5) 23 (2022), 2003--2048,
\href{https://arxiv.org/pdf/2003.06547}{arXiv:2003.06547v2},
especially \S\S2, 5--7, 9.
A literature search located this earlier conditional proof. We make no claim
to be the first to observe the consequence of the new upstream theorem.

\begin{corollary}[Coefficients in the CM field]
Let $K=\mathbb Q(\sqrt\Delta)$. The minimum support size stays exactly
$s(\Delta)$ when the coefficients range over $K$.
\end{corollary}
\begin{proof}
The ring class field theorem gives
$[K(\alpha):K]=h(\Delta)=[\mathbb Q(\alpha):\mathbb Q]$.
Consequently the natural map
$K\otimes_{\mathbb Q}\mathbb Q(\alpha)\to K(\alpha)$ is an isomorphism.
For any finite choice of exponents, independence of the corresponding powers
over $\mathbb Q$ therefore persists over $K$. Conversely a rational relation
is also a relation over $K$. The same argument works for any coefficient
field linearly disjoint from $\mathbb Q(\alpha)$ over $\mathbb Q$.
\end{proof}

\subsection{A density-one extension to any fixed number of terms}

For four or more terms, the preceding published trinomial theorem leaves the
minimum support undetermined. The same geometry of CM conjugates yields a bound
which can exceed any prescribed support size. The following calculation
uses the reduced forms of Bilu--Luca--Pizarro-Madariaga, \S2 and the proof
of Proposition 3.1(6), and the root-counting argument of their Lemma 12.3.
It holds independently of the repository's GRH theorem.

\begin{proposition}[Support from coprime factorizations]
Let $\omega(n)$ be the number of distinct prime divisors of $n$. Then
\[
\begin{aligned}
 s(-4n)&\ge 1+2^{\omega(n)-1} &&(n>1),\\
 s(-N)&\ge 1+\left\lceil 2^{\omega(N)-2}\right\rceil
 &&(N>3,\ N\equiv3\pmod4).
\end{aligned}
\]
In particular every order discriminant $-4n$ with at least three distinct
primes dividing $n$ excludes all four-term rational polynomials.
\end{proposition}
\begin{proof}
There are $2^{\omega(n)-1}$ unordered coprime factorizations $n=uv$,
$u<v$. Assign each full prime power to one of two factors and identify the
two orderings. The forms $[u,0,v]$ are distinct primitive reduced forms
of discriminant $-4n$. Their CM points are $i\sqrt{v/u}$. All give positive
conjugates, since for $z$ on the positive imaginary axis the formula
$j(z)=E_4(z)^3/(q\prod_{m\ge1}(1-q^m)^{24})$, $q=e^{2\pi iz}$, has
positive numerator and denominator by the Eisenstein series expansion.
Descartes' rule gives at most $k-1$ positive roots for a polynomial with
$k$ terms. This proves the even bound.

For $N=uv$ odd, again use coprime pairs $u<v$. They are opposite modulo
$4$. If $v>3u$, take $[u,u,(u+v)/4]$. If $v<3u$, take
$[(u+v)/4,(v-u)/2,(u+v)/4]$. Equality forces $(u,v)=(1,3)$, excluded.
Each triple is integral, reduced and has discriminant $-N$. For
primitivity, an odd common prime divisor would divide both $u$ and $v$,
and the middle coefficient is odd. Each pair can be recovered from its
triple, and the two types have disjoint shapes. The associated CM points
$z=(b+i\sqrt N)/(2a)$ have either $\Re z=1/2$ or $|z|=1$. In the first
case $-\bar z=z-1$ and in the second $-\bar z=-1/z$, so modular invariance
and the real Fourier coefficients of $j$ show that $j(z)$ is real. All
these conjugates are nonzero because the unique zero of $j$ on the
modular quotient has discriminant $-3$. Of $2^{\omega(N)-1}$ real
conjugates, at least half have the same sign. Apply Descartes' rule to
$P(X)$ or $P(-X)$ to finish.
\end{proof}

\begin{corollary}[Density for order discriminants]
Fix $k\ge1$. Among negative order discriminants, counted once each in
order of absolute value, the proportion with $s(\Delta)\le k$ tends to zero.
The same holds for equations with coefficients in the corresponding CM field.
\end{corollary}
\begin{proof}
For either branch of the proposition, bounded $s$ forces bounded $\omega$.
For clarity, fix $L$ and select the odd primes up to $y$. For
integers $m$ with natural-density weights, divisibility by these primes is a
family of independent indicators, by the Chinese remainder theorem.
Independence and the individual probabilities $1/p$ persist on
$m\equiv3\pmod4$. Their sum has mean $\mu_y=\sum_{2<p\le y}1/p$
and variance at most $\mu_y$. For $\mu_y>L$, Chebyshev's inequality bounds
the upper density of $\omega(m)\le L$ by $\mu_y/(\mu_y-L)^2$ in either
population. Divergence of the prime reciprocal sum lets $y$ tend to
infinity. Apply this to $m=n$ in the even case and to $m=N$ in the odd case.
The two cases exhaust the order discriminants up to finitely many exceptions.
\end{proof}

Two examples attain the bound. At $\Delta=-120$ the complete list of
reduced primitive forms is $[u,0,30/u]$ for $u=1,2,3,5$. At $\Delta=-840$
it is $[u,0,210/u]$ for $u=1,2,3,5,6,7,10,14$. The class polynomials have
degrees $4$ and $8$, so their support sizes are at most $5$ and $9$.
The even bound proves equality for the minimum over all polynomial
multiples, of arbitrary degree. The accompanying
\texttt{verify\_cm\_support.py} checks the reduced-form lists with integer
arithmetic. The density statement weights each discriminant once. Individual
odd prime discriminants can have arbitrarily large absolute value while
giving only one coprime factor pair, so the factor-pair bound alone gives no growing lower bound in that case.

\subsection{A uniform support bound from primes splitting in the CM field}

The zero-location theorem also gives a bound that grows with the discriminant
even when the discriminant has only one prime factor. We derive it by finding
many conjugates of different absolute sizes. The exponents and coefficients
of a potential sparse equation may be arbitrarily large.

\begin{lemma}[Separated conjugate sizes]
Let $z_1,\ldots,z_t$ be nonzero complex numbers whose moduli $r_i=|z_i|$
satisfy $r_{i+1}\ge16r_i$. For any integers $0\le e_1<\cdots<e_t$,
the matrix $(z_i^{e_j})_{i,j=1}^t$ is invertible. In particular a polynomial
vanishing on all these points has more than $t$ terms unless it is zero.
\end{lemma}
\begin{proof}
Multiply column $j$ by a positive number $c_j$, starting with $c_1=1$ and
choosing $c_{j+1}/c_j=(r_jr_{j+1})^{-(e_{j+1}-e_j)/2}$.
In row $i$ the ratio of successive absolute entries is
$(r_i/\sqrt{r_jr_{j+1}})^{e_{j+1}-e_j}$. It is at least $4$ up to the
diagonal ($j<i$) and at most $1/4$ subsequently ($j\ge i$). The sum of the
off-diagonal magnitudes in a row is at most $2\sum_{m\ge1}4^{-m}=2/3$
times its diagonal magnitude. Strict diagonal dominance implies
invertibility. To see this, take a putative nonzero kernel vector, choose a coordinate of
largest modulus and use the corresponding row. Column scaling is reversible.
For a polynomial with fewer than $t$ terms apply the same result to an
equal-sized subset of the points.
\end{proof}

\begin{theorem}[Uniform support growth]
With the Dirichlet GRH theorem as input, for every sufficiently large
absolute order discriminant $D=|\Delta|$,
\[
                  s(\Delta)\gg\frac{D^{1/4}}{\sqrt{\log D}} .
\]
The implied constant is absolute. Thus each fixed support size occurs at
only finitely many discriminants, including orders of arbitrary conductor.
\end{theorem}
\begin{proof}
Write $D=f^2|d_K|$ for $K=\mathbb Q(\sqrt\Delta)$, with $d_K$ its
fundamental discriminant. Set
$X=\lfloor D^{1/4}\sqrt{\log D}/10\rfloor$ and
$H=\lfloor\log X\rfloor$. Let $S$ consist of the odd primes in $(X/2,X]$
which split in $K$ and do not divide $f$.

\emph{Counting the split primes.}
The quadratic factorization $\zeta_K(s)=\zeta(s)L(s,\chi_{d_K})$
supplies RH for this specific Dedekind zeta function. The uniform prime
ideal theorem of Greni\'e and Molteni, Corollary 1.2, yields, for $y\ge100$,
\[
  |\psi_K(y)-y|
   \ll\sqrt y\big((1+\log y)\log |d_K|+1+\log^2y\big).
\]
Here $\psi_K$ assigns weight $\log N\mathfrak p$ to every power of a prime
ideal $\mathfrak p$. In degree two the contributions other than first powers
of split rational primes are $O(\sqrt y\log y+\log|d_K|)$.
Indeed, over each $p\le\sqrt y$ all prime-power terms have combined
weight $O(\log y)$, and the sum of $\log p$ at first powers of ramified
primes is at most $\log|d_K|$. A split $p$ contributes $2\log p$ at first
powers. Removing the primes dividing $f$ costs another $O(\log(2f))$.
Taking the difference at $y=X$ and $X/2$ therefore gives
$2\sum_{p\in S}\log p=X/2+o(X)$ uniformly in $\Delta$.
In particular $|S|\gg X/\log X$.
See L.~Greni\'e and G.~Molteni,
\href{https://arxiv.org/pdf/1312.4463}{\emph{Explicit versions of the
prime ideal theorem for Dedekind zeta functions under GRH}}.

\emph{Extracting spaced primes.}
For $1\le m\le X$ let $b_m=|S\cap(m,m+H]|$.
Then $\sum b_m=H|S|\gg X$, since each selected prime occurs for exactly
$H$ starts. The unconditional prime second-moment estimate of
Bazzanella, Languasco and Zaccagnini, Lemma 4 with $k=2$, states in our
range that $\sum_{m\le X}(\psi(m+H)-\psi(m))^2\ll X\log^2X$.
As $b_m\log(X/2)\le\psi(m+H)-\psi(m)$, this gives $\sum b_m^2\ll X$.
Cauchy--Schwarz shows that at least a constant times $X$ starts have
$b_m>0$. Take a maximal subset $S'$ of $S$ with pairwise distances at
least $H$. For every occupied start some $p\in S$ lies in $(m,m+H]$
and some $p'\in S'$ has $|p-p'|<H$. Each $p'$ accounts for at most
$3H+1$ such starts. Thus $t=|S'|\gg X/H$.
See D.~Bazzanella, A.~Languasco and A.~Zaccagnini,
\href{https://arxiv.org/pdf/0809.2967}{\emph{Prime numbers in logarithmic
intervals}}, \S3. Only the upper bound for the second moment is used.

\emph{Producing separated conjugates.}
For $p\in S'$, $(\Delta/p)=1$, $p\nmid\Delta$ and $4p^2<D$ for large
$D$. Choose a root of $b^2\equiv\Delta\pmod p$, impose the parity
$b\equiv\Delta\pmod2$, and take the representative $-p<b\le p$.
Then $c=(b^2+D)/(4p)$ is integral, $p<c$, and $[p,b,c]$ is a primitive
reduced form. Its conjugate $z_p$ satisfies
$\big||z_p|-\exp(\pi\sqrt D/p)\big|\le2079$ by the uniform estimate
in Bilu--Luca--Pizarro-Madariaga, \S2. Since $\sqrt D/X\to\infty$,
we have $\log|z_p|=\pi\sqrt D/p+o(1)$ uniformly. For $p<p'$ in
$S'$ the difference of these log-magnitudes is at least
\[
 \frac{\pi\sqrt D\,H}{X^2}-o(1)\longrightarrow 25\pi>\log16
\]
in the sense that the displayed lower bound tends to $25\pi$.
Order the $t$ conjugates by increasing modulus and apply the lemma. Every
annihilating polynomial has at least $t+1\gg X/H\gg D^{1/4}/\sqrt{\log D}$
terms.
\end{proof}

This proof of uniform growth uses two published estimates and an explicit
matrix argument. It does not determine the exceptional discriminants for
four terms. For an elementary variant avoiding the prime second moment,
take $X=D^{1/4}/10$ and use all eligible primes with spacing at least one.
The same proof then yields $s(\Delta)\gg D^{1/4}/\log D$.

There is also an unconditional qualitative comparison. For each fixed
$t$, choose $t$ relatively compact open disks in the modular curve with
positive hyperbolic area, such that every $j$-modulus in each disk is more
than sixteen times every $j$-modulus in the preceding disk. Such disks
exist by continuity near points $iy$ tending to the cusp, where
$j(iy)\sim e^{2\pi y}$. Equidistribution of full CM packets, including
nonmaximal orders, implies that for all sufficiently large $D$ the packet
meets each disk. Indeed one can integrate a nonnegative continuous bump of
positive integral supported in each disk. The lemma then gives
$s(\Delta)>t$. This independently yields qualitative finiteness for bounded
support. The split-prime argument gives the quantitative lower bound
$D^{1/4}/\sqrt{\log D}$. See L.~Clozel and E.~Ullmo,
\href{https://ullmo.ihes.fr/publications/clozel-ullmo.pdf}{\emph{\'Equidistribution
de points de Hecke}}, Theorem 2.4.

\begin{corollary}[A finite certificate with 7,958 conjugates]
For $\Delta=1-2^{61}$, every nonzero rational polynomial vanishing at
$j$ of discriminant $\Delta$ has at least $7{,}959$ terms, independently of
any zero-location theorem.
\end{corollary}
\begin{proof}
Let $D=2^{61}-1$ and $B=\lfloor D^{1/4}\rfloor=38{,}967$.
The accompanying integer sieve counts $7{,}958$ integers $1\le a\le B$
all of whose odd prime factors $p$ have $(\Delta/p)=1$.
For each such $a$, Hensel lifting at odd primes and the Chinese remainder
theorem show that $\Delta$ is a square modulo $4a$. The $2$-part follows
from $\Delta\equiv1\pmod{2^{61}}$. Every $a$ counted is coprime to
$\Delta$. Hence there is a primitive reduced form with leading coefficient
$a$ by Proposition 3.1(1) of Bilu--Luca--Pizarro-Madariaga.
For $a<a'\le B$, the exponents $\pi\sqrt D/a$ and $\pi\sqrt D/a'$
differ by at least $\pi\sqrt D/B^2\ge\pi$. Their exponentials are
both at least $e^{\pi B}$, so the absolute error $2079$ changes their
magnitudes by less than one percent. The actual ratio therefore exceeds
$(0.99/1.01)e^\pi>16$. Apply the separated-size lemma to all $7{,}958$
moduli. The computation needs neither a numerical value of $j$ nor the
coefficients of the class polynomial.
\end{proof}

\subsection{An explicit cutoff for equations with four terms}

\begin{theorem}[Four-term cutoff]
Using Dirichlet GRH as supplied by the upstream result,
$s(\Delta)\ge5$ for every negative order discriminant with
$D=|\Delta|\ge10^{11}$.
\end{theorem}
\begin{proof}
Put $R=\sqrt D$. Seek one odd split prime coprime to the conductor in
each of the three intervals $(uR,vR]$ in the table below. Write
\[
 E(y)=\sqrt y\left[\left(\frac{\log y}{2\pi}+2\right)\log D+
                  2\left(\frac{\log^2 y}{8\pi}+2\right)\right],
 \quad B(y)=2\sqrt y\log y+\log D+2\log2 .
\]
Put $E(0)=0$. Corollary 1.2 of Greni\'e--Molteni gives
$|\psi_K(y)-y|\le E(y)$ for $y\ge100$.
Twice the sum of $\log p$ for the desired primes in $(x,y]$ is at least
$y-x-E(x)-E(y)-B(y)$ when $x=0$ or $x\ge100$.
To see the removal cost, all prime ideal powers over the rational primes
$p\le\sqrt y$ contribute at most $2\sqrt y\log y$. At larger $p$,
inert primes make no contribution. Ramified primes cost at most
$\log|d_K|$, and removing the split primes at the conductor costs at most
$2\log f$. Their sum is $\log D$. The term $2\log 2$ permits the
explicit removal of $2$.

At $D=10^{11}$ the following are outward integer bounds, computed with
$316227<R<316228$ and $\pi>3.14$. The columns correspond to $x=uR$ and
$y=vR$. The accompanying \texttt{verify\_four\_term.py} checks them using
rational arithmetic, including the logarithm bounds.
\[
\begin{array}{c|r|r|r}
 (u,v] & \text{lower bound on }y-x & \text{upper bound on }E(x)+E(y)
      & \text{upper bound on }B(y) \\ \hline
 (0,0.046]       &14546&12138&2340\\
 (0.048,0.189]   &44588&38992&5405\\
 (0.230,0.500]   &85381&75082&9548
\end{array}
\]
All three margins are positive. They persist for larger $R$ after
normalizing by $R$. Indeed $E(cR)/R$ is a positive sum of constants
times $R^{-1/2}$ times products of at most two factors chosen from
$\log R,\log(cR)$. For $c\ge0.046$ these logarithms exceed $9$.
The derivative of the logarithm of each term with respect to $\log R$
is at most $-1/2+2/9<0$. The terms in $B(cR)/R$ decrease by the same
argument, including $R^{-1}\log D$. The normalized interval widths are
constant.

For each of the resulting primes $p_1,p_2,p_3$, Proposition 3.1(1) of
Bilu--Luca--Pizarro-Madariaga gives a reduced form of leading
coefficient $p_i$ since $p_i\nmid D$, $\Delta$ is a square modulo $4p_i$,
and $4p_i^2\le D$. Include the principal form of leading coefficient $1$.
For successive prime intervals ending at $v$ and beginning at $u'$, the
ratio of the associated larger modulus to the smaller one is at least
\[
       \frac{e^{\pi/v}-2079}{e^{\pi/u'}+2079}.
\]
Here $\pi(1/v-1/u')>2.84$ for both gaps and all four exponentials exceed
$e^{13}>207900$. Thus the ratio exceeds $(0.99/1.01)e^{2.84}>16$.
For the principal form versus $p_1\ge3$, the ratio is at least
$(e^{\pi R}-2079)/(e^{\pi R/3}+2079)>16$. The separated-size lemma
applies to the four resulting conjugates. This gives an explicit finite
range for the remaining four-term classification.
\end{proof}
